\documentclass[11pt]{article}
\usepackage[margin=1in]{geometry}
\usepackage[T1]{fontenc}
\usepackage[utf8]{inputenc}
\usepackage{lmodern}
\usepackage{microtype}
\usepackage{amsmath,amssymb,amsthm,mathtools}
\usepackage{enumitem}
\usepackage{hyperref}
\usepackage{xcolor}
\usepackage{listings}
\usepackage{booktabs}
\usepackage{longtable}
\usepackage{array}
\usepackage{titlesec}
\usepackage{csquotes}
\hypersetup{colorlinks=true,linkcolor=blue,citecolor=blue,urlcolor=blue}
\setlist[itemize]{leftmargin=1.5em}
\setlist[enumerate]{leftmargin=1.75em}
\lstset{basicstyle=\ttfamily\small,breaklines=true,columns=fullflexible,frame=single}

\title{Evidence-Carrying World Models, Typed Governance Semantics,\\ and Categorical Backends: A Conceptual Primer}
\author{}
\date{2026-03-09}

\begin{document}
\maketitle

\begin{abstract}
This note consolidates a family of conceptual ideas developed around world-model
calculus, evidence semantics, governance reasoning, and typed semantic backends.
Its aim is didactic rather than archival: to state the cleanest current picture,
separate what is already mathematically sharp from what remains conjectural, and
record a stable set of design commitments for future formalization work.

The central thesis is that a world-model calculus should be understood first as
an \emph{evidence evaluator}, not as a global proof theory.  The additive
fragment---where query extraction is a monoid homomorphism from posterior-state
revision into an evidence carrier---is mathematically distinguished and deserves
to remain a named canonical subclass.  More realistic semantics for provenance,
correlation, overlap, forgetting, and locality should be layered \emph{around}
that additive core rather than replacing it.  Modal, neighborhood,
first-order, higher-order, and probabilistic relational semantics can then be
hosted as semantic backends or extensional embeddings; governance and treaty
reasoning become typed application layers over the same substrate.

The note also proposes a categorical placement: externally, as an
\emph{evidence-valued institution}; internally, as an
\emph{evidence-valued hyperdoctrine} or doctrine of typed queries.  On the
syntax side, dependent type theory and categories with families provide the
natural candidate for the internal language of the trusted ontology/kernel
layer, while presheaf-topos infrastructure remains the right ambient semantic
home for the broader indexed semantics stack.
\end{abstract}

\tableofcontents

\section{Purpose and Scope}

This document is a conceptual primer for a research line centered on:
\begin{itemize}
  \item a \emph{world-model calculus} for evidence-carrying semantic evaluation,
  \item typed governance and treaty semantics built on top of that calculus,
  \item categorical transport layers (institution, hyperdoctrine, Beck--Chevalley),
  \item modal, neighborhood, first-order, higher-order, and probabilistic
        relational backends,
  \item and the possibility that a dependent type theory supplies the internal
        language of the trusted ontology layer.
\end{itemize}

The note is deliberately \emph{conceptual}.  It is not a changelog, not a
proof index, and not a full literature survey.  It aims instead to preserve the
strongest conceptual picture now available so that future development can resume
from a stable basis.

\section{The Core Idea: Evidence First, Views Second}

The starting point is a simple but important methodological separation:
\begin{quote}
\textbf{Probability is what one queries; evidence is what one carries.}
\end{quote}

The intended semantics is not that a world-model state directly \emph{is} a
probability distribution, a proof state, or a truth assignment.  Rather, a
world-model state carries enough information that queries can extract an
\emph{evidence object}, and then different operational views of that evidence
can be selected: strength, confidence, acceptance, interval bounds, posterior
parameters, and so on.

This separation matters for at least four reasons.
\begin{enumerate}
  \item It keeps raw support distinct from acceptance policy.
  \item It permits multiple views over the same carried evidence.
  \item It prevents confidence, probability, and actuality from being conflated.
  \item It makes semantic backends composable: multiple logics can target the
        same evidence layer while exposing different operational selectors.
\end{enumerate}

\subsection{The additive evaluator}

The cleanest core interface is:
\begin{lstlisting}
class WorldModel (State Query : Type*) [EvidenceType State] where
  evidence : State -> Query -> Evidence
  evidence_add : forall W1 W2 q,
    evidence (W1 + W2) q = evidence W1 q + evidence W2 q
\end{lstlisting}

The state type is equipped with revision/addition, and query extraction lands
in an additive evidence carrier.  This says that each query induces a monoid
homomorphism from the revision algebra of states into the evidence algebra.

This is the mathematically distinguished \emph{additive evidence ledger}
fragment.  It is not merely a convenient implementation choice.  It captures
exactly the semantics of independent sufficient-statistic aggregation in several
canonical families.

\section{Why the additive fragment is special}

The additive fragment should remain a named privileged subclass even if the
framework is generalized around it.

\subsection{Ledger before policy}

Exact additivity is the cleanest semantics of \emph{accumulating raw evidence}:
independent observations or independent posterior fragments add.  This records
how much support and counter-support have been gathered without yet deciding
what to trust, what to discount, what to forget, or what to accept.

Trust discounting, overlap correction, anti-double-counting, and forgetting are
all philosophically meaningful, but they are already \emph{epistemic policy}
choices.  They should therefore be layered around additive accumulation rather
than used to erase its identity.

\subsection{Canonical sufficient-statistic semantics}

The additive fragment is also singled out by conjugate Bayesian families.
For example:
\begin{itemize}
  \item Beta--Bernoulli: evidence counts are sufficient statistics.
  \item Dirichlet--multinomial: category counts are sufficient statistics.
  \item Normal--Gamma: count, sum, and sum-of-squares are sufficient statistics.
\end{itemize}

In each case, the sufficient statistics aggregate by addition.  This makes the
additive world-model fragment the canonical semantics of independent evidence
accumulation for a broad and useful class of empirical models.

\subsection{A universal-property viewpoint}

The mathematically strongest current story is not that additive world models are
the unique semantics of \emph{all} evidence, but that they are the unique
\emph{free additive extension} of atomic observations.

Given atomic observation contributions
\[
  a : \mathrm{Obs} \to \mathrm{Query} \to Ev,
\]
one obtains a unique additive extension
\[
  E : \mathrm{Multiset}\,\mathrm{Obs} \to \mathrm{Query} \to Ev
\]
satisfying:
\begin{align*}
E(\{o\},q) &= a(o,q), \\
E(\sigma_1 + \sigma_2,q) &= E(\sigma_1,q)+E(\sigma_2,q).
\end{align*}

The right long-term theorem target is a full classification result:
\begin{quote}
Every additive generic world model on multisets is uniquely determined by its
singleton observation surface.
\end{quote}

That would justify the additive layer as the canonical semantics of independent
atomic evidence.

\section{Why the generic shell is not vacuous}

A recurring source of confusion is the intentional triviality of the bare state
judgment.

\subsection{Admissible states versus proof theory}

The context-free world-model judgment should be read as:
\begin{quote}
Any posterior state is admissible input to the evaluator.
\end{quote}

This is \emph{not} the claim that any proposition is derivable.  Query
judgments remain deterministic because for fixed state and query the evidence is
fixed by the extractor.  So the shell behaves more like big-step evaluation than
like theoremhood.

The right rule of thumb is:
\begin{itemize}
  \item the bare shell is an evaluator host,
  \item trusted derivability begins only in provenance-restricted or
        context-indexed layers,
  \item and proof-like content enters through certified rewrite rules,
        consequence rules, and closure operators.
\end{itemize}

\subsection{Semantic nontriviality via endpoints}

The framework becomes nontrivial through instantiated backends and endpoint
correspondences.

For example, if one instantiates states as singleton pointed Kripke models or
singleton pointed neighborhood models, then query strength recovers ordinary
satisfaction truth values.  With soundness and completeness assumptions for the
external proof systems, implication-level proof-theoretic correspondences can be
transported into world-model strength inequalities.

Thus the generic shell is not a deep global proof system, but it is a faithful
semantic host for nontrivial proof-theoretic endpoints.

\section{Views are selectors, not semantics}

Strength, confidence, acceptance, and interval bounds should all be treated as
\emph{selectors} or \emph{views} of evidence, not as the semantics itself.

\subsection{No universal ``better evidence'' order}

Raw evidence often comes with an obvious algebraic order, but that order need
not coincide with ``better evidence'' for all purposes.  In particular, an order
that treats both positive and negative mass as monotone can make confidence
increase while support quality decreases.

A more coherent approach is to define \emph{view-induced preorders}.  For
example, with a strength context $ctx$ and confidence parameter $\kappa$ one may
use:
\[
  e_1 \preceq_{ctx,\kappa} e_2
  \iff
  \mathrm{strengthWith}(ctx,e_1) \le \mathrm{strengthWith}(ctx,e_2)
  \;\wedge\;
  \mathrm{confidence}(\kappa,e_1) \le \mathrm{confidence}(\kappa,e_2).
\]

This keeps kernel algebra and acceptance policy conceptually distinct.

\subsection{Actuality as evidence, not metaphysics}

A closely related clarification concerns actuality.  A ``really exists'' or
\texttt{rexist} tag should not be read as metaphysical certainty.  In an
evidence-based world-model calculus, the correct reading is:
\begin{quote}
actuality is an \emph{event-occurrence evidence channel}, and categorical
``actuality for governance'' is a downstream acceptance policy.
\end{quote}

This yields a clean two-layer design:
\begin{itemize}
  \item \texttt{eventualityEvidence / eventEvidence}
  \item \texttt{acceptedOccurrence / eventAccepted}
\end{itemize}

The second is a policy-relative categorical layer; the first is the semantic
one.

\section{Subsumption, compliance, and treaty semantics}

A major application layer developed around these ideas is a typed governance and
treaty semantics.

\subsection{Event subsumption}

Practical governance does not only compare exact events.  It often asks whether
an abstract obligation is satisfied by a more specific concrete event.

The clean solution is a subsumption relation on eventualities:
\begin{itemize}
  \item same predicate,
  \item same polarity,
  \item and every thematic role specified in the abstract event is matched by
        the concrete one.
\end{itemize}

This is the right formal replacement for brittle exact matching when checking
whether an observed event fulfills an abstract clause.

\subsection{Treaty kernels}

On top of event subsumption, a minimal treaty kernel becomes natural:
\begin{itemize}
  \item treaty clauses with obligor, obligee, modality, event pattern,
        optional deadline, and optional repair clause,
  \item timestamped treaty traces,
  \item separate predicates for \emph{obligation fulfillment} and
        \emph{prohibition violation},
  \item and later policy-gated accepted traces.
\end{itemize}

The conceptual point is that compliance should not be post-hoc rule checking
only.  Treaty semantics should influence planning and repair.  This is a useful
minimal formal layer for an AI-to-AI treaty economy.

\section{Modal and logical backends}

The world-model calculus is broad enough to host several kinds of semantics.

\subsection{Kripke semantics}

A Kripke backend can be obtained by taking states to be multisets of pointed
models and letting query evidence count supporting versus refuting points.
Singleton states then recover ordinary pointed-model truth.

This gives a clean adequacy bridge:
\begin{quote}
singleton query strength is ordinary modal truth;
multisets yield evidential ensemble semantics over modal hypotheses.
\end{quote}

\subsection{Neighborhood semantics}

Neighborhood semantics is even more native to evidence-style reasoning.  It
supports weaker and non-normal modalities, and it aligns naturally with
families of evidence sets rather than only relational accessibility.

In the current conceptual picture:
\begin{itemize}
  \item Kripke semantics is the clearest classical adequacy witness,
  \item neighborhood semantics is the better long-term evidence-native modal
        backend.
\end{itemize}

\subsection{First-order and higher-order semantics}

The same evaluator pattern can host set-theoretic FOL and Henkin-style HOL:
\begin{itemize}
  \item states become multisets of pointed structures and assignments,
  \item queries become formulas in context,
  \item singleton adequacy recovers ordinary model satisfaction,
  \item and implication-level closure wrappers can transport proof-theoretic
        consequence into world-model strength inequalities.
\end{itemize}

This suggests that the world-model calculus is not tied to modal logic.  It is
better viewed as a semantics carrier for multiple logical families.

\section{Probabilistic relational semantics: ProbLog and MLNs}

ProbLog and Markov Logic Networks fit naturally at the extensional semantic
level.

\subsection{Weighted-world compilation}

Both frameworks induce weighted possible-world semantics.  One can compile them
into a weighted-world world-model state and define evidence by splitting total
world weight into support and refutation mass for a query.

Then query strength is the induced marginal probability of the query.  In this
sense:
\begin{quote}
ProbLog and MLNs can be \emph{subsumed extensionally} by the world-model
calculus.
\end{quote}

\subsection{What is and is not preserved}

This does \emph{not} mean that the current additive state-revision operation is
the native program-combination semantics of ProbLog or MLNs.  Their intensional
composition laws are global and non-additive in important ways.

So the correct claim is:
\begin{itemize}
  \item exact query-level semantic subsumption: yes,
  \item exact native syntactic composition under additive revision: no.
\end{itemize}

The right architecture is therefore a weaker evaluator core with the additive
world-model retained as the canonical independent-evidence specialization.

\section{Categorical placement}

The correct categorical home is not a single topos claim and not the claim that
this calculus is itself the syntactic category of a theory.

\subsection{Externally: evidence-valued institution}

Externally, the cleanest formulation is as an \emph{evidence-valued
institution}:
\begin{itemize}
  \item a category of signatures or contexts,
  \item sentences/queries over each signature,
  \item models/world-model states over each signature,
  \item and evidence-valued satisfaction stable under change of notation.
\end{itemize}

This captures the idea that the calculus is a semantic target for many logical
front ends.

\subsection{Internally: evidence-valued hyperdoctrine}

Internally, the strongest picture is hyperdoctrinal:
\begin{itemize}
  \item typed queries over a categorical base,
  \item reindexing along substitutions or morphisms,
  \item existential and universal transport,
  \item Beck--Chevalley coherence.
\end{itemize}

This is exactly the right setting for first-order and typed query semantics.

\subsection{Ambient home: presheaf-topos flavored semantics}

Presheaf topoi are the right \emph{ambient} semantic home for much of the stack:
strong enough to support internal intuitionistic logic, predicate fibrations,
change of base, and Yoneda-style context analysis, but generic enough not to
prematurely overdetermine the ontology.

They are not, by themselves, the final semantics of the trusted dependent type
kernel.

\section{Dependent type theory and the trusted ontology layer}

The right relation between the world-model calculus and dependent type theory is
subtle but promising.

\subsection{What DTT can provide}

A dependent type theory with contexts, substitution, $\Pi$/$\Sigma$-types,
identity types, and universes provides the natural \emph{syntax side} or
internal language of the trusted ontology layer.

Categorically, this is better described via:
\begin{itemize}
  \item categories with families (CwFs),
  \item locally cartesian closed categories (LCCCs),
  \item or related contextual/category-with-attributes structures,
\end{itemize}
not merely by saying ``topos.''

\subsection{What DTT does not automatically provide}

Dependent type theory does not by itself provide:
\begin{itemize}
  \item an evidence algebra,
  \item posterior-state revision,
  \item provenance weighting,
  \item locality/gluing policies,
  \item or experiment-selection semantics.
\end{itemize}

So the best long-term slogan is:
\begin{quote}
Dependent type theory can supply the internal language and context machinery of
the trusted ontology/kernel layer; the world-model calculus supplies the
empirical/evidential semantics layered on top of it.
\end{quote}

\subsection{The right theorem direction}

The strongest future theorem pack here would show that:
\begin{enumerate}
  \item a dependent type theory generates the relevant syntactic/contextual
        category,
  \item a world-model hyperdoctrine interprets those contexts and substitutions,
  \item substitution is respected,
  \item and key type formers are interpreted by the expected semantic
        operations.
\end{enumerate}

That would make the world-model calculus a semantic doctrine into which the
trusted syntax layer is interpreted.

\section{Why sheaves and subtopoi belong below the world-model layer}

Locality, observers, time slices, and patchwise ontology growth are not reasons
to rewrite the world-model surface.  They are reasons to enrich its semantic
backends.

\subsection{Presheaf first, sheaf later}

A presheaf backend gives local data plus restriction maps.  A sheaf backend adds
a gluing law: compatible local sections assemble into global sections.

This becomes useful when one has:
\begin{itemize}
  \item observer-local semantics,
  \item distributed partial knowledge,
  \item context-sensitive concepts,
  \item local ontology patches that should sometimes glue.
\end{itemize}

\subsection{Why crisp gluing is too strict}

For AGI-style growth, exact equality on overlaps is often too strong.  Local
patches may only be partially compatible, evidence-weightedly compatible, or
paraconsistently compatible.

So the correct long-term target is not merely ordinary sheaf gluing but:
\begin{quote}
\emph{evidence-valued or paraconsistent gluing over a locality backend.}
\end{quote}

That allows local-to-global organization without destroying useful partial
syntheses.

\section{Fixpoint closure and operational reasoning}

The world-model calculus already has the right seed for operational reasoning:
fixpoint closure.

\subsection{Why fixpoints matter}

Once a one-step consequence operator is available, one can define:
\begin{itemize}
  \item immediate consequence,
  \item least closure,
  \item bounded cascades,
  \item eventual discovery on finite spaces,
  \item and later greatest/guarded variants.
\end{itemize}

This is the right semantic story for iterative reasoning, recursive closure,
and stabilization.

\subsection{The $\mu$-to-$\rho$ connection}

If an operational backend such as a $\rho$-calculus is used for execution, the
best architecture is:
\begin{itemize}
  \item keep world-model/fixpoint semantics as the denotational spec,
  \item prove a $\mu$-style query/fixpoint embedding into the operational
        process backend,
  \item show that one-step consequence and eventual stabilization correspond to
        operational evolution.
\end{itemize}

This is much stronger than replacing fixpoint semantics with process execution.
It preserves a semantic anchor.

\section{Provenance, acceptance, and trusted traces}

For practical AGI use, provenance and policy-gated acceptance must become
first-class.

\subsection{Provenance-first semantics}

A strong default story is:
\begin{itemize}
  \item raw observations/events are wrapped with provenance,
  \item evidence is extracted from those assessed items,
  \item acceptance policies decide what enters the trusted trace,
  \item reasoning and treaty compliance then run over accepted traces.
\end{itemize}

This avoids the false dichotomy between ``mere evidence'' and ``already
accepted fact.''

\subsection{Forgetfulness should be explicit}

Forgetting is not default revision.  It should be a separate operation with its
own laws, such as idempotence and invariance for queries outside the forgotten
scope.

This is crucial for:
\begin{itemize}
  \item privacy and redaction,
  \item stale-evidence expiry,
  \item theory contraction,
  \item and explicit support-scope management.
\end{itemize}

\section{Experiment design and formal science}

One of the most important future directions is to connect the world-model
calculus to the design and evaluation of experiments.

\subsection{Why this matters}

A semantics of evidence is not yet a semantics of science.  Science also needs:
\begin{itemize}
  \item experiment selection,
  \item value of information,
  \item comparison of information channels,
  \item and policy-relative optimality.
\end{itemize}

The right long-term mathematical interfaces here likely come from categorical
probability and experiment comparison (e.g. Blackwell-style orderings and
Markov-category inspired transport), while the world-model calculus supplies the
semantic state/evidence/query layer into which those experiment morphisms
report.

\section{A recommended hierarchy of world-model structures}

The best long-term formulation seems to be a hierarchy, not a single monolithic
class.

\begin{enumerate}
  \item \textbf{WorldModelCore}: evidence extractor only.
  \item \textbf{Generic merge layer}: possibly weak, provenance-aware merge.
  \item \textbf{Overlap/correlation layer}: overlap correction or compatibility.
  \item \textbf{AdditiveWorldModel}: exact additivity under independent
        accumulation.
  \item \textbf{Forgetting layer}: explicit contraction/redaction.
  \item \textbf{Locality backend}: presheaf first, sheaf/subtopos later.
  \item \textbf{Closure/fixpoint layer}: consequence and stabilization.
  \item \textbf{Typed semantic transport layer}: institution + hyperdoctrine.
  \item \textbf{Trusted ontology/kernel bridge}: dependent type theory or other
        certified syntax side.
\end{enumerate}

The existing additive world-model should remain the named canonical subclass,
not be weakened in place.

\section{A practical theorem roadmap}

The following theorems would most strongly stabilize the theory.

\subsection{For the additive fragment}
\begin{enumerate}
  \item \textbf{Classification theorem}: every additive generic world model on
        multisets is uniquely induced by its singleton observation surface.
  \item \textbf{Functoriality}: additive extension is natural under additive
        homomorphisms of evidence carriers.
  \item \textbf{No-go results}: additive plus idempotent merge plus nontriviality
        are incompatible; additive plus exact inverse forgetting is incompatible.
  \item \textbf{Representation theorem}: additive sufficient-statistic semantics
        induce additive world models.
\end{enumerate}

\subsection{For non-additive extensions}
\begin{enumerate}
  \item \textbf{Overlap-corrected merge}: additive recovery under explicit
        independence.
  \item \textbf{Trust action laws}: source discounts act functorially on evidence.
  \item \textbf{Forgetting invariance}: queries outside forgotten scope are
        unaffected.
  \item \textbf{Evidence-valued gluing}: local compatibility induces global
        assembly under policy.
\end{enumerate}

\subsection{For the syntax bridge}
\begin{enumerate}
  \item interpretation of substitution/reindexing,
  \item interpretation of key dependent type formers,
  \item reduction soundness into world-model consequence,
  \item and ideally an initiality/universality theorem for the syntax-to-WM
        interpretation.
\end{enumerate}

\section{Working assumptions for future development}

For future discussion and implementation work, the following assumptions should
be treated as the current clean baseline.

\begin{enumerate}
  \item The world-model calculus is primarily a semantics/evaluation interface,
        not a global proof system.
  \item The additive fragment is mathematically distinguished and should remain
        explicit.
  \item Strength, confidence, acceptance, and intervals are views of evidence,
        not the evidence itself.
  \item ``Actuality'' should be modeled as event-occurrence evidence together
        with policy-relative acceptance, not as metaphysical certainty.
  \item Governance reasoning benefits from event subsumption and treaty kernels,
        not only exact event matching.
  \item Kripke and neighborhood semantics are backends or endpoints, not rivals
        to the core evaluator.
  \item Dependent type theory is a plausible syntax/internal-language side of the
        trusted ontology layer, but not a replacement for evidence semantics.
  \item Locality, gluing, and observer semantics belong below the world-model
        interface as backend refinements.
  \item Fixpoint closure should remain the denotational spec for iterative
        reasoning even if operational backends are added.
  \item Experimental design and value-of-information semantics are natural next
        layers for turning evidence semantics into a formal science engine.
\end{enumerate}

\section{Conclusion}

The most robust conceptual picture is now clear.

A world-model calculus should be treated as an evidence-carrying semantic
substrate with a mathematically distinguished additive fragment, a typed
categorical transport layer, and multiple logical/probabilistic backends.
Dependent type theory should be viewed as a candidate syntax/internal-language
side of the trusted ontology layer, not as a replacement for empirical
semantics.  Provenance, acceptance, forgetting, locality, and experiment
design should be layered on top or below the core in their proper places rather
than collapsed into a single overloaded notion of ``truth.''

The additive fragment is special, but not absolute.  It is the right semantics
of independent evidence accumulation, not the unique semantics of all epistemic
policy.  The correct long-term architecture therefore preserves that fragment as
a canonical center while generalizing around it in a disciplined hierarchy.

\begin{thebibliography}{99}

\bibitem{GoguenBurstallInstitution}
Joseph Goguen and Rod Burstall.
\newblock \emph{Institutions: Abstract model theory for specification and
programming}.
\newblock Journal of the ACM, 39(1):95--146, 1992.

\bibitem{LawvereHyperdoctrine}
F.~William Lawvere.
\newblock \emph{Adjointness in foundations}.
\newblock Dialectica, 23(3--4):281--296, 1969.

\bibitem{DybjerCWF}
Peter Dybjer.
\newblock \emph{Internal type theory}.
\newblock In \emph{Types for Proofs and Programs}, 1995.

\bibitem{KupkePattinsonCoalgebraic}
Clemens Kupke and Dirk Pattinson.
\newblock \emph{Coalgebraic semantics of modal logics: an overview}.
\newblock Theoretical Computer Science, 2009.

\bibitem{vanBenthemEvidence}
Johan van Benthem, David Fern{\'a}ndez-Duque, and Eric Pacuit.
\newblock \emph{Evidence and plausibility in neighborhood structures}.
\newblock Annals of Pure and Applied Logic, 165(1):106--133, 2014.

\bibitem{RichardsonDomingosMLN}
Matthew Richardson and Pedro Domingos.
\newblock \emph{Markov logic networks}.
\newblock Machine Learning, 62(1--2):107--136, 2006.

\bibitem{DeRaedtProbLog}
Luc De Raedt, Angelika Kimmig, and Hannu Toivonen.
\newblock \emph{ProbLog: A probabilistic Prolog and its application in link
 discovery}.
\newblock IJCAI, 2007.

\bibitem{VanBenthemEtAl}
Johan van Benthem, David Fern{\'a}ndez-Duque, and Eric Pacuit.
\newblock \emph{Evidence logic and evidence dynamics}.
\newblock Synthese, 196:2567--2593, 2019.

\bibitem{Walley}
Peter Walley.
\newblock \emph{Statistical Reasoning with Imprecise Probabilities}.
\newblock Chapman and Hall, 1991.

\bibitem{MarkovCategories}
Tobias Fritz.
\newblock \emph{A synthetic approach to Markov kernels, conditional
independence, and theorems on sufficient statistics}.
\newblock Advances in Mathematics, 370:107239, 2020.

\bibitem{BlackwellCats}
Tobias Fritz and coauthors.
\newblock \emph{Blackwell comparison in Markov categories}.
\newblock Theoretical Computer Science, 2023.

\end{thebibliography}

\end{document}
